\section{Background and Related Work}
\subsection{Task Scope and Security Evaluation}
Security benchmarks differ in their inputs and success criteria. PrimeVul studies code-language-model vulnerability detection while addressing label accuracy, duplication, chronological splits, and evaluation realism~\cite{primevul}. CORRECT evaluates context-rich vulnerable/patched program pairs and the rationales accompanying classification~\cite{correct}. These concerns motivate caution when interpreting coverage of a small set of annotated locations as complete security analysis.

VulnGym aligns advisories with vulnerable repository snapshots and annotates EPs, COs, and connecting traces~\cite{vulngym}. Its evaluation includes end-to-end entry prediction and oracle-conditioned subtasks. We use its verified entries but define a narrower acquisition comparison: the public description and dataset-supplied entry title are given, explicit EP/CO fields are withheld, and the model returns an undifferentiated list of candidate file/line intervals. We score file coverage and tolerance-based interval hits. These scores do not reproduce VulnGym's complete entry/trace evaluation and are not numerically interchangeable with its benchmark results.

CyberGym uses executable proof-of-concept reproduction from descriptions and codebases~\cite{cybergym}, providing a different form of validation. Project-scale empirical work compares specialized LLM detectors, general-purpose agents, and traditional static analyzers, examining both detection outcomes and costs~\cite{projectscale}. IRIS combines LLM-inferred taint specifications and contextual analysis with static analysis~\cite{iris}. Our configurations perform source acquisition and location prediction without a static-analysis engine, rationale adjudicator, or exploit validator; none of these systems is treated as an experimentally reproduced baseline.

\subsection{Context Acquisition and Interaction}
Repository context can be selected before inference or acquired through feedback. RepoCoder combines retrieval and generation iteratively for repository-level code completion~\cite{zhang-etal-2023-repocoder}. Our single-pass RAG shares the use of repository retrieval but has a different task and interaction pattern. It combines sparse and dense rankings using reciprocal rank fusion~\cite{rrf}. Agent chooses read-only actions after observing earlier results; Pack supplies a task-independent lexical prefix.

\subsection{Directly Related Evaluation Settings}
VLoc Bench~\cite{vloc} studies file localization from a generic CWE description, withholding advisory and patch hints. It pairs vulnerable and patched snapshots, scores patch-derived files with file F1, and measures negative-state behavior through a common read-only terminal. It also classifies budget exhaustion, premature termination, wrong-file submissions, and partial localization in its failure analysis. Our dataset-supplied titles, EP/CO annotations, and coverage-only scores define a different task with no patched negatives; output-reliability analysis itself already has this precedent.

Agent Retrieval Bench~\cite{agentretrieval} evaluates files needed next in coding workflows, using workflow-derived queries, ranked recall, budgeted context yield, no-gold controls, logged acquisition, and a controlled context-seed pilot. Its token-budget exposure metric is adjacent to our target-source measurement, but its targets and tokenizer differ. We instead jointly record acquisition-ledger tokens, cumulative provider input, and final localization output for three advisory-conditioned configurations.

VulContextBench~\cite{vulcontext} gives an introducing diff and candidate repository snapshot, and scores role-tagged context at file, block, and line granularity. It separately measures viewed and declared context under mini-SWE-agent. Thus the distinction between acquired and returned evidence already has direct security precedents. Our incremental analysis combines shared nonblank eligibility with annotation-level source/output outcomes, the contribution of invalid outputs to paired contrasts, prediction extent, and separate acquisition/input ledgers in a same-model configuration comparison.
% Generated by generate_literature_comparison.py from verified primary-source notes.
\begin{table*}[t]
\centering\small
\caption{Nearby evaluation settings. The three cited benchmark papers are preprints; their scores are not compared numerically with this study.}\label{tab:nearby}
\begin{tabular}{@{}>{\raggedright\arraybackslash}p{0.17\textwidth}>{\raggedright\arraybackslash}p{0.24\textwidth}>{\raggedright\arraybackslash}p{0.25\textwidth}>{\raggedright\arraybackslash}p{0.25\textwidth}@{}}
\toprule
Setting & Given input & Targets and evaluation & Source/scaffold and budget boundary \\
\midrule
VLoc Bench~\cite{vloc} & Generic CWE description; advisory and fix withheld & Patch-derived vulnerable files; file F1, patched negatives, and non-submission/failure categories & Common read-only terminal; command and generation-turn limits \\
Agent Retrieval Bench~\cite{agentretrieval} & Pull-request, review, failure-trace, and anchored-edit signals & Next-context files; ranked recall, budgeted yield, no-gold controls & Frozen base corpus; regex-token packing budgets; logged agents and a seed pilot \\
VulContextBench~\cite{vulcontext} & Introducing diff and read-only candidate snapshot & Role-tagged files/blocks/lines; viewed and declared precision/recall & Common mini-SWE-agent; numbered inspected evidence and final submission \\
This study & Dataset entry title (may name target file), truncated public description, and pinned snapshot & Withheld EP/CO files and intervals; coverage only, no negative snapshots & One model, three configurations; acquisition ledger A and cumulative input B \\
\bottomrule
\end{tabular}
\end{table*}


Table~\ref{tab:nearby} summarizes the input, target, negative-control, and budget differences. These preprints inform study positioning; we neither compare their absolute scores with ours nor claim that evidence/output separation itself is new. The comparison concerns complete implemented configurations, including presentation, interface, history replay, and answer validation.
